% Ad Atticum 14.2 (ad-atticum-14-02)
% Source: https://raw.githubusercontent.com/PerseusDL/canonical-latinLit/master/data/phi0474/phi057/phi0474.phi057.perseus-lat1.xml
% Fetched by scripts/fetch_latin.py

% Dateline (Perseus): Scr. in suburbano Mati vi Id. Apr. a. 710 (44).

\ciceroLetterOpener{CICERO ATTICO salutem}

\ciceroSection{1} duas a te accepi epistulas heri. ex priore theatrum Publiliumque cognovi, bona signa consentientis multitudinis. plausus vero L. Cassio datus etiam facetus mihi quidem visus est.

\ciceroSection{2} altera epistula de Madaro scripta, apud quem nullum † φαλάκωμα , ut putas. processit enim, sed minus diutius. sermone eius† sum retentus. quod autem ad te scripseram obscure fortasse, id eius modi est. aiebat Caesarem secum, quo tempore Sesti rogatu veni ad eum, cum exspectarem sedens, dixisse, ego nunc tam sim stultus ut hunc ipsum facilem hominem putem mihi esse amicum quom tam diu sedens meum commodum exspectet? habes igitur φαλάκρωμα inimicissimum oti, id est Bruti.

\ciceroSection{4} in Tusculanum hodie, Lanuvi cras, inde Asturae cogitabam. Piliae paratum est hospitium, sed vellem Atticam. verum tibi ignosco. quarum utrique salutem.
